% Do not change document class, margins, fonts, etc.
\documentclass[a4paper,oneside,bibliography=totoc]{scrbook}

% some useful packages (add more as needed)
\usepackage[utf8]{inputenc}
\usepackage{graphicx}
\usepackage{latexsym}
\usepackage{amsmath}
\usepackage{amssymb}
\usepackage{tabularx}
\usepackage{booktabs}
\usepackage{algorithm} % you can modify the algorithm style to your liking
\counterwithin{algorithm}{chapter} % so that algorithms have chapter numbers as well
\usepackage{algorithmic}
\usepackage{csquotes}
\renewcommand{\algorithmiccomment}[1]{\hfill\textit{// #1}}
\usepackage[usenames,dvipsnames]{xcolor}
\usepackage[colorlinks,citecolor=Green]{hyperref} % you may change/remove the colors
% lipsum package removed: not available on this cluster's TeX install and only used by template placeholder text below

% chicago citation style
\usepackage{natbib}
\bibliographystyle{plainnat} % chicagoa.bst is unavailable on this cluster's TeX install; plainnat ships with natbib and works with the authoryear \setcitestyle below; switch back to chicagoa if your submission environment has it
\setcitestyle{authoryear,round,semicolon,aysep={},yysep={,}} \let\cite\citep

% example enviroments (add more as needed)
\newtheorem{definition}{Definition} \newtheorem{proposition}{Proposition}

%remove this: 
% \textcolor{red}{TODO: }
\usepackage{xcolor}
\usepackage[htt]{hyphenat}

\begin{document}

\frontmatter \subject{Master Thesis} % change to appropriate type
\title{Design and Evaluation of Efficient Semantic Company Retrieval at Scale}
\author{Marmee Pandya\\
  (matriculation number 1963521)} \date{\today}
\publishers{{\small Submitted to}\\
  Chair of Enterprise Data Analysis\\
   Dr.\ Ralph Peeters\\
  University of Mannheim\\}
\maketitle

\chapter{Abstract}

Your thesis must contain an abstract. A good reference for thesis writing
is~\citet{zobel2014}; we highly recommend that you study this or a similar book
during your studies. He writes the following about the abstract:

\blockcquote{zobel2004}{%
  An abstract is typically a single paragraph of about 50 to 200 words. The
  function of an abstract is to allow readers to judge whether or not the paper
  is of relevance to them. It should therefore be a concise summary of the
  paper's aims, scope, and conclusions. There is no space for unnecessary text;
  an abstract should be kept to as few words as possible while remaining clear
  and informative. Irrelevancies, such as minor details or a description of the
  structure of the paper, are inappropriate, as are acronyms, abbreviations, and
  mathematics. Sentences such as "We review relevant literature" should be
  omitted.

  The more specific an abstract is, the more interesting it is likely to be.
  Instead of writing "space requirements can be significantly reduced", write
  "space requirements can be reduced by 60\%". Instead of writing "we have a new
  inversion algorithm", write "we have a new inversion algorithm, based on
  move­to-front lists".

  Many scientists browse research papers outside their area of expertise. You
  should not assume that all likely readers will be specialists in the topic of
  their paper-abstracts should be self-contained and written for as broad a
  readership as possible. Only in rare circumstances should an abstract cite
  another paper (for example, when one paper consists entirely of analysis of
  results in another), in which case the reference should be given in full, not
  as a citation to the bibliography.}

% table of contents
\begingroup%
\hypersetup{hidelinks}% disable link color in TOC only
\tableofcontents%
\endgroup

% none of the things below is needed, but you may add them if you feel that they
% are helpful for your work \listofalgorithms \listoffigures \listoftables
% \listtheorems{definition} \listtheorems{proposition}


% okay, start new numbering ... here is where it really starts
\mainmatter
% here the original content of my thesis starts 

\chapter{Project Overview}
\label{ch:projectOverview}

\texttt{Thesis Goal:} Build and evaluate an efficient embedding-based semantic retrieval system for large-scale company data that outperforms the current production search system.

\noindent\texttt{Core Question:} Can embedding-based semantic retrieval combined with approximate nearest neighbor search efficiently generate accurate top-k candidate company sets from large-scale datasets while maintaining low query latency? \newline

\noindent\textbf{Current Todo List:}
- Literature Review

% TODO (when writing the Introduction chapter): this Project Overview can
% become the basis for the Introduction. Its final paragraph should be a
% roadmap paragraph explaining what each subsequent chapter covers
% (Chapter 2 covers X, Chapter 3 covers Y, and so on).

\chapter{ISTARI.AI and its Data}
\label{ch:istariAndData}

\section{About ISTARI.AI}

This thesis is conducted in collaboration with ISTARI.AI\protect\footnote{\url{https://docs.istari.ai/goi/goi_start_page/}}, a data intelligence company that specialises in providing high-quality firmographic data at scale. Istari's core product is a curated, organisation-level dataset that covers organisations worldwide, enabling use cases such as market analysis, company search, and lead generation. The central research challenge addressed in this thesis is improving the quality of semantic company retrieval, arises directly from Istari's operational needs.

\section{The Global Organisation Index (GOI)}

The primary dataset used in this thesis is Istari's \textbf{Global Organisation Index} (GOI), a curated collection of approximately 20 million active organisations spanning 232 countries and territories. The GOI serves as the canonical source of firmographic data within Istari's data infrastructure, covering organisation identity, location, industry classification, size, etc.

\section{Data Pipeline and Methodology}

The GOI is constructed through a multi-stage validation pipeline that deliberately prioritises data quality over volume. The pipeline proceeds as follows:

\begin{enumerate}
    \item \textbf{Registry Collection.} National registers are systematically queried worldwide, including commercial registers, company registers, and association registers. These primary sources are enriched with additional open data sources, yielding an initial pool of approximately 400 million organisations.

    \item \textbf{Domain Attribution.} Each organisation is assessed for attribution to a clearly identifiable web domain. Approximately 10\% of the initial pool, around 40 million organisations can be matched to a domain. The remainder are excluded as they are either no longer active, were never truly operational (e.g., pure holding structures), or lack a web presence.

    \item \textbf{Activity Verification.} The 40 million domain-attributed organisations are further filtered by verifying that their domains are still actively operated. This final step reduces the dataset to approximately 20 million organisations that are demonstrably active.
\end{enumerate}

This pipeline ensures that every record in the GOI is both web domain-verified and confirmed active, a key differentiator from traditional firmographic databases, which typically retain large numbers of inactive or dormant records.

\subsection{Schema and Key Fields}

Each record in the GOI contains a set of core identity fields, geographic fields, and classification fields. The fields most relevant to this thesis are summarised in Table~\ref{tab:goi_schema}.

\begin{table}[h]
\centering
\caption{Key fields of the GOI schema used in this thesis.}
\label{tab:goi_schema}
\begin{tabular}{lll}
\toprule
\textbf{Field} & \textbf{Type} & \textbf{Description} \\
\midrule
\texttt{name}              & STRING  & Organisation name \\
\texttt{domain}            & STRING  & Organisation website domain (unique identifier) \\
\texttt{summary}           & STRING  & AI-generated textual summary of the organisation \\
\texttt{keywords}          & LIST    & Keywords \\
\texttt{employee\_class}            & STRING  & Employee count bracket \\
\texttt{country}            & STRING  & Country \\
\texttt{country\_code}            & STRING  &  ISO Country Code \\
\texttt{state / state\_code}            & STRING  & State or province \\
\texttt{region / region\_code}            & STRING  & Region \\
\texttt{district / district\_code}            & STRING  & District \\
\texttt{municipality / muncipality\_code}            & STRING  & Municipality \\
\texttt{address}            & STRING  & Full address \\
\texttt{latitude}            & STRING  & Latitude coordinate \\
\texttt{longitude}            & STRING  & Longitude coordinate \\
\texttt{nace\_code}        & STRING  & NACE industry classification code \\
\texttt{organization\_type}& STRING  & Company, Startup, Academic, Public, or Other \\
\texttt{organization\_size}& STRING  & Micro, Small, Medium-sized, or Large enterprise \\
\texttt{country}           & STRING  & Country of registration \\
\bottomrule
\end{tabular}
\end{table}

Of particular importance for this thesis is the \texttt{summary} field. This field contains an AI-generated natural language description of each organisation, synthesised from its web presence and available structured data. It is the most semantically rich and consistently populated field in the dataset, with effectively zero missing values across the corpus used for experiments, and therefore serves as the primary input for all retrieval models evaluated in this thesis.

\subsection{Industry Classification: NACE Codes}

Industries in the GOI are classified using the \textbf{NACE} taxonomy \protect\footnote{\url{https://ec.europa.eu/eurostat/web/nace}} (\textit{Nomenclature statistique des activités économiques dans la Communauté européenne}), the European Union's standard statistical classification of economic activities. The GOI covers 22 NACE sectors. Table~\ref{tab:nace_top} lists the ten most represented sectors in the dataset.

\begin{table}[h]
\centering
\caption{Top 10 industries in the GOI by organisation count.}
\label{tab:nace_top}
\begin{tabular}{llrr}
\toprule
\textbf{NACE} & \textbf{Industry} & \textbf{Count} & \textbf{Share (\%)} \\
\midrule
N & Professional, scientific \& technical activities & 3{,}093{,}879 & 18.03 \\
G & Wholesale \& retail trade                        & 2{,}240{,}191 & 13.05 \\
C & Manufacturing                                    & 2{,}006{,}435 & 11.69 \\
R & Human health \& social work                      & 1{,}440{,}340 &  8.39 \\
K & Telecom, IT \& computing                         & 1{,}269{,}131 &  7.39 \\
F & Construction                                     & 1{,}173{,}226 &  6.84 \\
I & Accommodation \& food service                    & 1{,}126{,}594 &  6.56 \\
S & Arts, sports \& recreation                       & 1{,}101{,}592 &  6.42 \\
Q & Education                                        &   584{,}430 &  3.41 \\
J & Publishing, broadcasting \& content              &   554{,}985 &  3.23 \\
\bottomrule
\end{tabular}
\end{table}

\subsection{Organisation Size Classification}

Organisation size in the GOI is derived from employee and revenue signals and bucketed into four categories following EU SME definitions, as shown in Table~\ref{tab:size_dist}.

\begin{table}[h]
\centering
\caption{Organisation size distribution in the GOI.}
\label{tab:size_dist}
\begin{tabular}{llrr}
\toprule
\textbf{Size} & \textbf{Employee Range} & \textbf{Count} & \textbf{Share (\%)} \\
\midrule
Micro           & 0--9   & 7{,}341{,}838 & 42.78 \\
Small           & 10--49 & 7{,}494{,}409 & 43.67 \\
Medium-sized    & 50--249 & 1{,}501{,}655 &  8.75 \\
Large enterprise & 250+  &   824{,}126 &  4.80 \\
\bottomrule
\end{tabular}
\end{table}

The distribution is heavily skewed towards smaller organisations: over 86\% of the GOI consists of micro and small enterprises. This has implications for retrieval, as queries targeting large enterprises represent a smaller and more distinctive subset of the corpus.

\subsection{Geographic Coverage}

The GOI covers 232 countries and territories. Table~\ref{tab:geo_top} lists the top five countries by volume.

\begin{table}[h]
\centering
\caption{Top 5 countries by organisation count in the GOI.}
\label{tab:geo_top}
\begin{tabular}{lr}
\toprule
\textbf{Country} & \textbf{Organisations} \\
\midrule
United States   & 3{,}626{,}794 \\
Germany         & 1{,}828{,}181 \\
United Kingdom  & 1{,}013{,}026 \\
Netherlands     &   778{,}416 \\
Italy           &   625{,}754 \\
\bottomrule
\end{tabular}
\end{table}

It is worth noting that approximately 2.7 million organisations in the GOI do not have a standardised administrative region. These records are intentionally retained as they remain valuable for non-geographic analyses such as industry or domain-level retrieval, the primary focus of this thesis.

\subsection{Evaluation Dataset}

The full 20-million-company GOI corpus is not directly used for retrieval in this thesis. Instead, Istari provided an evaluation dataset consisting of 101 real-world search queries along with each query's top-1\,000 production search results, yielding 101\,000 (query, company) pairs in total. Where the same domain appears linked to multiple addresses, the GOI deduplicates by domain and retains the record with the highest employee count, treating it as the organisation's headquarters; after this deduplication, the evaluation corpus covers \textbf{98\,716 unique companies}. All retrieval experiments in Chapter~\ref{ch:experiments} are conducted over this corpus. The exact pseudo-relevance threshold used to define ``relevant'' companies for evaluation (and how it changes across experiments) is described in Sections~\ref{subsec:eval_protocol} and~\ref{sec:protocol_note}.

Testing retrieval quality at the full 20-million-company scale was not pursued in the current phase of this thesis, since doing so would require a substantially larger data export from Istari than the query-level top-1\,000 results used here; this is discussed as an open question in the concluding chapter.

\chapter{Literature}
\label{ch:related_literature}

This chapter reviews the retrieval paradigms, indexing techniques, and
evaluation practices that this thesis builds on and compares against.
Section~\ref{sec:lit_sparse_dense} covers sparse and dense retrieval,
the two foundational paradigms for the baselines in
Chapter~\ref{ch:experiments}. Section~\ref{sec:lit_ann} covers
approximate nearest neighbour search, the technique that makes dense
retrieval practical at scale. Section~\ref{sec:lit_late_interaction}
covers late-interaction retrieval and retrieval-specific pretraining,
the two alternative paradigms evaluated in
Sections~\ref{sec:simlm}--\ref{sec:colbert}.
Section~\ref{sec:lit_hybrid} covers hybrid retrieval, cross-encoder
reranking, and learning to rank, the techniques underlying the hybrid
pipelines in Sections~\ref{sec:hybrid_fullpool}--\ref{sec:learned_fusion}.
Section~\ref{sec:lit_evaluation} covers evaluation practice in
information retrieval when human relevance judgements are unavailable,
directly motivating the pseudo-relevance protocol used throughout this
thesis (Section~\ref{subsec:eval_protocol}).

% ────────────────────────────────────────────────────────────
\section{Sparse and Dense Retrieval}
\label{sec:lit_sparse_dense}
% ────────────────────────────────────────────────────────────

Text retrieval systems fall into two broad families. \textbf{Sparse}
retrieval represents queries and documents as high-dimensional vectors
of term weights and scores matches by term overlap. The dominant sparse
method is BM25 \citep{robertson2009probabilistic}, a probabilistic
scoring function that weighs each query term by its frequency in a
document and its rarity across the whole collection. BM25 is fast,
requires no training, and remains a strong baseline decades after its
introduction, but it cannot match a query and a document that describe
the same concept using different words: this is the vocabulary mismatch
problem, directly observed for this thesis's task in
Section~\ref{subsec:bm25_qual}.

\textbf{Dense} retrieval instead represents queries and documents as
fixed-size vectors learned by a neural encoder, and scores matches by
vector similarity rather than literal term overlap. Dense Passage
Retrieval (DPR) \citep{karpukhin2020dense} established the now-standard
bi-encoder architecture for open-domain retrieval: a query encoder and a
document (or, in this thesis, company) encoder, trained so that relevant
pairs are close in vector space, encoding each side independently so
that documents can be embedded once, offline, and reused across every
future query. Sentence-BERT \citep{reimers2019sentence} showed that the
same bi-encoder idea, applied to general-purpose sentence embeddings
rather than a retrieval-specific training objective, produces embeddings
useful for semantic similarity and retrieval tasks broadly, and underlies
the general-purpose sentence-transformer models (including MiniLM,
Section~\ref{subsec:minilm}) used throughout this thesis. More recent
embedding models, such as the BGE family \citep{chen2024bge}, combine
this bi-encoder architecture with large-scale contrastive pretraining and
hard-negative mining, and represent the current state of the art for
general-purpose dense retrieval, forming the basis of several baselines
evaluated in Sections~\ref{subsec:bge} and~\ref{sec:extended_baselines}.

% ────────────────────────────────────────────────────────────
\section{Approximate Nearest Neighbour Search}
\label{sec:lit_ann}
% ────────────────────────────────────────────────────────────

Dense retrieval requires finding the nearest vectors to a query vector
among the full document collection. Exact search costs time linear in
the number of documents, which becomes impractical at the scale of
millions of documents relevant to this thesis's target deployment
(Section~\ref{ch:istariAndData}). FAISS \citep{johnson2017billion}
provides GPU-accelerated implementations of both exact and approximate
search, including inverted-file indexes (IVF), which cluster the
embedding space and search only the clusters nearest the query, trading
a small amount of recall for a large reduction in query time. Hierarchical
Navigable Small World graphs (HNSW) \citep{malkov2020efficient} take a
different approach: they build a multi-layer proximity graph over the
document vectors and greedily traverse it at query time, again trading
exactness for speed. Both IVF and HNSW are evaluated directly in this
thesis's ANN scaling experiment (Section~\ref{sec:ann_scaling}), which
tests whether the accuracy--speed trade-off these methods make remains
favourable at simulated production scale.

% ────────────────────────────────────────────────────────────
\section{Late Interaction and Retrieval-Specific Pretraining}
\label{sec:lit_late_interaction}
% ────────────────────────────────────────────────────────────

A bi-encoder compresses an entire document into one fixed-size vector,
which can discard fine-grained information relevant to a specific query.
ColBERT \citep{khattab2020colbert} addresses this with \emph{late
interaction}: query and document are each represented as a set of
per-token vectors, and relevance is scored by summing, for every query
token, the maximum similarity to any document token (MaxSim), so no
single vector has to summarise the whole document in advance. This adds
compute and storage cost relative to a single-vector bi-encoder, since
retrieval now requires comparing token sets rather than single vectors.
ColBERTv2 \citep{santhanam2022colbertv2} introduces residual compression
of the per-token vectors and the PLAID engine for efficient late-interaction
search at scale, substantially reducing this cost; this thesis instead
implements MaxSim scoring directly (Section~\ref{sec:colbert}), trading
the PLAID engine's efficiency for an exact, dependency-light
implementation suited to the scale of the evaluation corpus used here.

A separate line of work improves dense retrieval not through
architecture but through pretraining. SimLM \citep{wang2022simlm}
pretrains a bi-encoder with a representation-bottleneck objective before
any retrieval-specific fine-tuning, explicitly designed to produce
embeddings well suited to downstream passage retrieval. This thesis
evaluates whether that retrieval-specific pretraining transfers to a
domain (company profiles) quite different from the passage-retrieval
benchmarks it was designed for (Section~\ref{sec:simlm}).

% ────────────────────────────────────────────────────────────
\section{Hybrid Retrieval, Reranking, and Learning to Rank}
\label{sec:lit_hybrid}
% ────────────────────────────────────────────────────────────

Sparse and dense retrieval, and different dense models, make different
errors: a hybrid system that combines multiple retrieval signals can
recover from a single method's individual weaknesses. Reciprocal Rank
Fusion (RRF) \citep{cormack2009reciprocal} is a simple, training-free way
to combine several ranked lists: each document's fused score is the sum
of $1/(k + \mathrm{rank})$ across the lists it appears in, requiring only
that a document be ranked, not that its raw scores be comparable across
methods. RRF underlies the hybrid baseline in Section~\ref{sec:rrf}.

A complementary technique is \emph{reranking}: a first-stage retriever
returns a candidate set cheaply, and a more expensive model reorders it.
Cross-encoder rerankers, such as the \texttt{bge-reranker-v2-m3}
\citep{baaireranker2024} and \texttt{mxbai-rerank-large-v1}
\citep{mixedbread2024mxbai} models used in
Sections~\ref{sec:hybrid_fullpool} and~\ref{sec:reranker_ablation}, feed
the query and each candidate into the same transformer jointly, rather
than encoding them independently as a bi-encoder does. This lets the
model attend directly between query and candidate tokens, typically
producing more accurate relevance scores than a bi-encoder, at the cost
of requiring a full forward pass per candidate rather than a single fast
vector comparison, so cross-encoders are used to rerank a small candidate
set rather than to search a full collection.

More generally, combining multiple retrieval signals into a single
relevance score can be framed as \emph{learning to rank}
\citep{liu2009learning}: instead of a fixed, hand-designed combination
rule such as RRF, a model is trained directly on labelled query-document
pairs to learn how to weight and combine whatever features are available
(similarity scores, ranks, or other signals) into a final ranking. This
is the approach taken by the learned fusion ranker in
Section~\ref{sec:learned_fusion}, implemented with the logistic
regression and gradient-boosted tree models provided by scikit-learn
\citep{pedregosa2011scikit}.

% ────────────────────────────────────────────────────────────
\section{Evaluation Without Human Relevance Judgements}
\label{sec:lit_evaluation}
% ────────────────────────────────────────────────────────────

Standard information retrieval evaluation compares each system's
ranking against a set of human relevance judgements for a fixed query
set, historically constructed through pooling: the top results from
several systems are merged and judged by human assessors, as established
by the TREC evaluation methodology \citep{voorhees1999trec}. No such
human-annotated relevance judgements exist for the company search task
addressed in this thesis. In their absence, this thesis instead adopts a
\emph{pseudo-relevance} protocol (Section~\ref{subsec:eval_protocol}):
a company is treated as relevant if it appears in the existing production
system's own results for that query. This is a common practice when
human judgements are unavailable, and the BEIR benchmark
\citep{thakur2021beir} similarly demonstrates that retrieval methods can
be compared consistently across many tasks and domains without
task-specific human annotation, provided the resulting scores are
interpreted as agreement with a reference system rather than as absolute
measures of relevance, a limitation this thesis states explicitly
wherever pseudo-relevance labels are used.

% ============================================================
\chapter{Experiments}
\label{ch:experiments}
% ============================================================

This chapter describes the experimental setup, baseline retrieval systems,
evaluation methodology, and results obtained for the company search task
over Istari's Global Organisation Index~(GOI).
Section~\ref{sec:exp_setup} introduces the dataset and evaluation protocol.
Sections~\ref{sec:baseline_sparse}--\ref{sec:baseline_dense} describe each
retrieval baseline in turn.
Section~\ref{sec:exp_results} presents and analyses the results across all
methods, and Section~\ref{sec:exp_runtime} provides a runtime comparison.
Section~\ref{sec:protocol_note} documents a change in evaluation protocol
adopted for all subsequent experiments.
Sections~\ref{sec:extended_baselines}--\ref{sec:oracle_union} extend the
baseline comparison, explore query expansion and hybrid fusion, validate
the pseudo-relevance labels against an LLM judge, compare against Istari's
production APIs, evaluate approximate nearest-neighbour search at scale,
and analyse the sources of retrieval error.
Sections~\ref{sec:simlm}--\ref{sec:colbert} evaluate two further retrieval
paradigms (a retrieval-pretrained encoder and late-interaction retrieval).
Sections~\ref{sec:hybrid_fullpool}--\ref{sec:learned_fusion} develop a
full-pool reranking pipeline, isolate the reranker as a bottleneck, and
close part of the remaining gap with a learned fusion ranker trained on
this task's own data, the best-performing method in this thesis.

% ────────────────────────────────────────────────────────────
\section{Experimental Setup}
\label{sec:exp_setup}
% ────────────────────────────────────────────────────────────

\subsection{Dataset}
\label{subsec:dataset}

The evaluation corpus is derived from Istari's production search logs.
A set of 101 representative queries was provided by Istari, covering a
wide range of company search intents, including industry-specific queries
(e.g.\ \textit{``pharmaceutical companies in the US''}), geographic
queries (e.g.\ \textit{``tech startups in Berlin''}), and
functional queries (e.g.\ \textit{``B2B SaaS companies''}).
For each query, the production system returned its top-1\,000 ranked
results, yielding 101\,000 query--company pairs in total.
After de-duplication across queries, the corpus consists of
\textbf{98\,716 unique companies}, each represented by a domain name,
a company name, and an AI-generated natural-language summary of the
company's business activity.

The summary field was selected as the primary retrieval field after
an exploratory data analysis of all available GOI attributes.
It was found to be fully populated across the corpus (0\% missing values)
and to contain the richest semantic content for matching against
free-text queries, compared to structured fields such as industry codes
or founding year.

\subsection{Evaluation Protocol}
\label{subsec:eval_protocol}

\paragraph{Pseudo-relevance labels.}
No human-annotated relevance judgements are available for the Istari
company search task.
Following standard practice in information retrieval research when
ground-truth labels are absent~\citep{voorhees1999trec}, this thesis
adopts a \textit{pseudo-relevance} assumption: a company is considered
\textbf{relevant} for a given query if and only if it appears in the
\textbf{top-100 results} of Istari's production search system for
that query.
This yields a relevance set of at most 100 companies per query,
representing ${\sim}0.1\%$ of the total corpus.

This assumption has an important implication for the evaluation: the
metrics measure the degree to which each retrieval system agrees with
the production system, not whether it recovers objectively relevant
companies.
Consequently, a retrieval system that is strictly better than the
production system may appear to underperform because it surfaces
companies the production system missed.
This limitation is acknowledged throughout the analysis.

\paragraph{Metrics.}
Three standard information retrieval metrics are reported at cut-offs
$k \in \{10, 50, 100, 1000\}$:

\begin{itemize}
    \item \textbf{Precision@$k$}: the fraction of the top-$k$ retrieved
    companies that are relevant.
    This measures user-facing quality: what a user sees on the first
    page of results.

    \item \textbf{Recall@$k$}: the fraction of all relevant companies
    that appear in the top-$k$ results.
    This measures coverage, important for downstream re-ranking
    pipelines that operate on a retrieved candidate set.
    \textbf{Recall is the primary metric in this thesis}, following the
    priority set by Istari for this task: whether the system finds the
    relevant companies matters more than the exact order they are
    ranked in. Section~\ref{sec:protocol_note} explains how Precision@$k$
    and Recall@$k$ relate to one another and should be read together.

    \item \textbf{NDCG@$k$} (Normalised Discounted Cumulative Gain):
    a ranking-aware metric that rewards systems for placing relevant
    companies higher in the list, using a logarithmic discount. NDCG@$k$
    is reported as a secondary, supporting metric: it reflects the
    experience of a user who interacts primarily with the top results,
    which is a reasonable concern in practice even though it is not the
    priority metric for this thesis.
\end{itemize}

All metrics are averaged across all 101 queries.
Precision and Recall are always reported at the same cut-off values
so that the precision--recall trade-off at each depth is visible.

\subsection{Computational Environment}
\label{subsec:compute}

All experiments are conducted on the University of Mannheim HPC cluster.
Dense model encoding is performed on an NVIDIA A100 GPU with 80\,GB
HBM2 memory.
BM25 retrieval runs on CPU.
Query latency measurements are performed on the same hardware and reported
as the mean over 20 benchmark queries to obtain stable estimates.
Latency is decomposed into three components (local query encoding,
FAISS index search, and external network time) so that self-hosted and
API-based methods are compared on equal terms.

% ────────────────────────────────────────────────────────────
\section{Sparse Retrieval Baseline: BM25}
\label{sec:baseline_sparse}
% ────────────────────────────────────────────────────────────

\subsection{Implementation}
\label{subsec:bm25_impl}

The BM25 baseline is implemented using the \texttt{rank\_bm25} Python
library~\citep{robertson2009probabilistic}, which provides the Okapi
BM25 variant with default hyperparameters $k_1 = 1.5$ and $b = 0.75$.
All 98\,716 company summaries are tokenised using the NLTK punkt
tokeniser~\citep{nltk} with lower-casing applied, and the resulting
token lists are used to construct the BM25 inverted index.
At query time, each query is tokenised identically and scored against
the index.
The top-1\,000 companies by BM25 score are retrieved for each query
and stored for evaluation.
Index construction takes 28.7 seconds and requires no GPU.

\subsection{Qualitative Analysis}
\label{subsec:bm25_qual}

Manual inspection of BM25 results reveals a consistent failure mode
arising from the \textit{vocabulary mismatch problem}: BM25 retrieves
companies that contain query keywords in their description, regardless
of whether those companies operate in the queried domain.

A representative example is the query \textit{``software companies''}.
The top BM25 results are predominantly investment banks and venture
capital firms whose summaries frequently mention the word \textit{software}
in the context of describing their investment portfolio
(e.g.\ \textit{``we invest in software, fintech, and SaaS businesses''}).
These companies are not software companies: they fund them.
BM25 cannot distinguish between a company \textit{operating} in a
domain and a company \textit{discussing} that domain.
This limitation provides a clear and concrete motivation for
semantic retrieval methods that operate on the meaning of text
rather than on token overlap.

% ────────────────────────────────────────────────────────────
\section{Dense Retrieval Baselines}
\label{sec:baseline_dense}
% ────────────────────────────────────────────────────────────

Three dense retrieval baselines are evaluated, covering a range of
model sizes, computational requirements, and deployment constraints.
All three follow the same two-stage pipeline introduced by
\citet{karpukhin2020dense}: (i) encode all company summaries offline
into a dense vector index using FAISS~\citep{johnson2017billion},
and (ii) at query time, encode the query into a vector and retrieve
the top-1\,000 nearest neighbours from the index.

\subsection{all-MiniLM-L6-v2 (MiniLM)}
\label{subsec:minilm}

\texttt{all-MiniLM-L6-v2} is a lightweight sentence-transformer
model~\citep{reimers2019sentence} with 22 million parameters,
producing 384-dimensional embeddings.
It is distilled from a larger BERT model and fine-tuned on sentence
similarity tasks.
All 98\,716 company summaries are encoded in a single batch pass,
taking 0.6 minutes on the A100 GPU.
A \texttt{faiss.IndexFlatL2} index (exact Euclidean search) is
constructed over the resulting embeddings in 0.2 seconds.
MiniLM does not require embedding normalisation.

MiniLM is included as the smallest and most computationally efficient
dense baseline, providing a lower bound on dense retrieval quality
and establishing whether a minimal dense model is sufficient to
outperform BM25.

\subsection{BGE-large-en-v1.5 (BGE)}
\label{subsec:bge}

\texttt{BAAI/bge-large-en-v1.5} is a 335-million parameter BERT-large
model developed by the Beijing Academy of Artificial
Intelligence~\citep{chen2024bge}, producing 1\,024-dimensional embeddings.
It is specifically optimised for dense retrieval through hard negative
mining during fine-tuning, a training strategy in which the model
is exposed to query--document pairs that are superficially similar
but semantically distinct, forcing it to learn fine-grained semantic
distinctions.

Encoding all 98\,716 company summaries takes 7.0 minutes on the A100 GPU
at a batch size of 256.
BGE requires L2 normalisation of all embeddings
(\texttt{normalize\_embeddings=True}), as it is trained under the
cosine similarity objective.
Accordingly, a \texttt{faiss.IndexFlatIP} index (inner product,
equivalent to cosine similarity for unit-normalised vectors) is used
rather than \texttt{IndexFlatL2}.

At query time, each query is encoded independently and normalised
before being passed to the FAISS index.
The query encoding step takes 16.6\,ms on the A100 GPU, and the
subsequent FAISS search takes 27.8\,ms, for a total query latency
of 44.4\,ms.

\subsection{OpenAI text-embedding-3-small (OpenAI)}
\label{subsec:openai}

\texttt{text-embedding-3-small} is a commercial embedding model
provided by OpenAI via their REST API.
It produces 1\,536-dimensional embeddings.
Unlike the open-source models above, the architecture and training
procedure are proprietary.
Encoding all 98\,716 company summaries is performed via batched API
calls (batch size 500) and takes 9.7 minutes, of which the majority
is network latency rather than model computation.
The resulting embeddings are unit-normalised by the API, and a
\texttt{faiss.IndexFlatIP} index is constructed locally.
The total cost of encoding the full corpus is approximately \$0.15
at the published pricing of \$0.02 per million tokens.

This model is included as the commercial upper-bound baseline,
following a suggestion from the industry supervisor that a
high-quality API-based embedding model be evaluated alongside
open-source alternatives.

\paragraph{Latency note.}
Query latency for the OpenAI baseline comprises two components:
a network round-trip to the OpenAI API servers (211.7\,ms, measured
over 20 queries) and a local FAISS search (42.9\,ms).
The network component is not a property of the model itself;
it reflects internet latency between the University of Mannheim cluster
and OpenAI's infrastructure, and should not be interpreted as a
measure of model speed.
For a fair comparison, the local FAISS search time (42.9\,ms)
is comparable to BGE's search time (27.8\,ms), confirming that
both methods perform equivalently at the index search stage.
The substantive difference is that OpenAI requires an external
API call for query encoding, whereas BGE encodes locally on GPU
in 16.6\,ms.

% ────────────────────────────────────────────────────────────
\section{Results}
\label{sec:exp_results}
% ────────────────────────────────────────────────────────────

\subsection{Retrieval Quality}
\label{subsec:quality}

Recall is the primary metric reported in this thesis, following the
priority set by Istari for this task: coverage of relevant companies
matters more than precise ranking of the first page of results.
Table~\ref{tab:main_results} presents Recall@$k$ for all four methods;
NDCG and Precision are discussed alongside it in the paragraphs below,
where individual values are relevant, rather than tabulated in full.

\begin{table}[ht]
\centering
\caption{Recall@$k$ of all four baseline systems averaged over
         101 queries. Bold indicates the best value in each column.}
\label{tab:main_results}
\begin{tabular}{lcccc}
\toprule
\textbf{Method} & \textbf{Recall@10} & \textbf{Recall@50} & \textbf{Recall@100} & \textbf{Recall@1000} \\
\midrule
BM25   & 0.016 & 0.068 & 0.128 & 0.682 \\
MiniLM & 0.019 & 0.069 & 0.122 & 0.653 \\
BGE    & \textbf{0.022} & \textbf{0.087} & \textbf{0.159} & 0.796 \\
OpenAI & \textbf{0.022} & 0.080 & 0.149 & \textbf{0.803} \\
\bottomrule
\end{tabular}
\end{table}

\paragraph{BM25 vs dense retrieval.}
All three dense methods outperform BM25 at low cut-offs.
At $k=10$, BGE and OpenAI achieve an NDCG of 0.227, representing a
\textbf{37\% improvement} over BM25's 0.166.
Precision@10 improves from 0.156 (BM25) to 0.216 (BGE), a 38\%
relative gain.
This confirms the hypothesis that semantic retrieval better captures
the intent of natural-language company search queries, particularly
for cases where relevant companies use vocabulary different from the
query terms.

At large cut-offs ($k = 1000$), the gap narrows.
BM25 achieves a Recall@1000 of 0.682, compared to BGE's 0.796 and
OpenAI's 0.803.
This pattern is consistent with findings in the BEIR
benchmark~\citep{thakur2021beir}: dense models are stronger at
ranking precision, while keyword-based methods maintain competitive
coverage at large retrieval depths due to their broad lexical matching.

\paragraph{BGE vs OpenAI.}
BGE and OpenAI are \textbf{effectively equivalent} across all metrics
and cut-offs.
The largest observed difference is 0.013 NDCG points (at $k=50$),
which is within the range of query-level variance and should not be
interpreted as a meaningful quality advantage.
This finding has an important practical implication: a free,
self-hosted model (BGE) achieves the same retrieval quality as a
commercial API (OpenAI) on this task.

\paragraph{MiniLM.}
MiniLM consistently underperforms BGE and OpenAI at all cut-offs,
and is only marginally better than BM25 at large $k$
(Recall@1000: 0.653 vs 0.682 for BM25).
It is also slower than BGE in production (query latency 80.9\,ms vs
44.4\,ms), as discussed in Section~\ref{sec:exp_runtime}.
MiniLM is therefore dominated by BGE on both quality and speed,
and there is no operating point at which it represents the optimal
choice.

\paragraph{Precision and the low-prevalence setting.}
Absolute precision values are low across all methods
(Precision@10 ranging from 0.156 to 0.224).
This is expected: relevant companies represent at most 100 out of
98\,716 corpus entries, approximately 0.1\% of the corpus.
In such a low-prevalence setting, even a perfect retrieval system
would show low absolute precision at large $k$.
What matters is the relative improvement between methods: BGE and
OpenAI achieve 38--43\% higher Precision@10 than BM25, which is a
meaningful gain.

\paragraph{The NDCG dip at $k = 100$.}
An artifact of the evaluation design is visible in
Table~\ref{tab:main_results}: NDCG values decrease from $k=10$
to $k=100$ before rising sharply at $k=1000$ for all methods.
This occurs because the pseudo-relevance labels are defined as the
production top-100.
At exactly $k=100$, NDCG measures whether all 100 relevant companies
are ranked perfectly within the top-100 retrieved results, a maximally
strict condition.
At $k=1000$, there is substantially more room to accommodate the relevant
set, and the metric rises accordingly.
This artifact would not appear with human-annotated relevance labels
that include varying degrees of relevance beyond a hard top-100 cut-off.

% ────────────────────────────────────────────────────────────
\section{Runtime Analysis}
\label{sec:exp_runtime}
% ────────────────────────────────────────────────────────────

Table~\ref{tab:latency} presents the decomposed query latency for
all four methods.
Latency is broken into three components to enable a fair comparison
between self-hosted models and the API-based OpenAI baseline.

\begin{table}[ht]
\centering
\caption{Query latency decomposed into local encoding, FAISS search,
         and external network time, averaged over 20 benchmark queries
         on an NVIDIA A100 GPU. Network time reflects internet latency
         to OpenAI servers and is not a property of the model.}
\label{tab:latency}
\setlength{\tabcolsep}{8pt}
\begin{tabular}{lrrrrl}
\toprule
\textbf{Method}
& \textbf{Encode (ms)}
& \textbf{Search (ms)}
& \textbf{Network (ms)}
& \textbf{Total (ms)}
& \textbf{Self-hostable} \\
\midrule
BM25   &   0.1 & 104.0 &   0.0 & 104.1 & Yes \\
MiniLM &  68.9 &  12.0 &   0.0 &  80.9 & Yes \\
BGE    &  16.6 &  27.8 &   0.0 &  \textbf{44.4} & Yes \\
OpenAI & \textit{server} &  42.9 & 211.7 & 254.6 & No \\
\bottomrule
\end{tabular}
\end{table}

\paragraph{BGE is the fastest method overall.}
BGE achieves a total query latency of 44.4\,ms, faster than BM25
(104.1\,ms), MiniLM (80.9\,ms), and OpenAI (254.6\,ms).
This result is initially counterintuitive: BGE is a 335-million
parameter neural network, whereas BM25 is a simple keyword scoring
function.
The explanation lies in hardware utilisation.
BGE encodes the query vector on the A100 GPU (16.6\,ms), a
highly parallelised operation that is extremely efficient for dense
linear algebra.
BM25, by contrast, computes token scores on CPU over an inverted
index (104.0\,ms), which is sequential and memory-bound.
This finding has a direct implication for production deployment:
a GPU-equipped system running BGE delivers both better retrieval
quality and lower query latency than a CPU-based BM25 system.

\paragraph{MiniLM encodes slowly despite being a small model.}
MiniLM's 68.9\,ms encoding time is slower than BGE's 16.6\,ms,
despite having 15 times fewer parameters.
This is because MiniLM encodes on CPU in these experiments, while
BGE encodes on GPU.
Hardware assignment dominates over model size for single-query
encoding latency.

\paragraph{OpenAI's latency is dominated by network time.}
Of OpenAI's total 254.6\,ms, only 42.9\,ms is FAISS index search;
comparable to BGE's 27.8\,ms search time.
The remaining 211.7\,ms is the HTTP round-trip to OpenAI's servers.
This latency is not a property of the embedding model; it reflects
internet connectivity between the university cluster and OpenAI's
infrastructure, and would vary across deployment environments.
For a production system at Istari, API-based encoding would introduce
an irreducible network dependency and per-query cost, whereas BGE
could be self-hosted with predictable latency.

\paragraph{Summary.}
Table~\ref{tab:summary} summarises the key trade-offs across all four
methods.

\begin{table}[ht]
\centering
\caption{Summary comparison of all baseline retrieval methods.}
\label{tab:summary}
\setlength{\tabcolsep}{6pt}
\begin{tabular}{lllllll}
\toprule
\textbf{Method}
& \textbf{NDCG@10}
& \textbf{Prec@10}
& \textbf{Recall@1k}
& \textbf{Latency}
& \textbf{Cost}
& \textbf{GPU req.} \\
\midrule
BM25   & 0.166 & 0.156 & 0.682 & 104\,ms & Free    & No \\
MiniLM & 0.192 & 0.186 & 0.653 & 81\,ms  & Free    & No \\
BGE    & \textbf{0.227} & 0.216 & 0.796 & \textbf{44\,ms} & Free & Yes (enc.) \\
OpenAI & \textbf{0.227} & \textbf{0.224} & \textbf{0.803} & 255\,ms & \$0.02/M tok. & No \\
\bottomrule
\end{tabular}
\end{table}

BGE emerges as the recommended baseline for subsequent experiments:
it matches OpenAI's retrieval quality, is the fastest method in
practice, is free, and can be fully self-hosted without API dependency.
The results motivate the next experimental phase, in which a
pre-filtering stage and hybrid retrieval system are built on top of
the BGE baseline to further improve precision at low cut-offs
and overall system efficiency.

% ────────────────────────────────────────────────────────────
\section{Evaluation Protocol Change for Follow-Up Experiments}
\label{sec:protocol_note}
% ────────────────────────────────────────────────────────────

All experiments described in the remainder of this chapter
(Sections~\ref{sec:extended_baselines}--\ref{sec:oracle_union}) use a
\textbf{relaxed pseudo-relevance definition}: a company counts as
relevant if it appears anywhere in production's \textbf{top-1\,000}
results for a query, rather than the top-100 threshold used in
Section~\ref{subsec:eval_protocol} and Table~\ref{tab:main_results}.
The relaxed threshold was adopted so that reranking, fusion, and
oracle-ceiling experiments would have enough relevant companies per
query to produce meaningful comparisons over large candidate pools.
As a direct consequence, the absolute NDCG/Precision/Recall values
reported from here on are on a different (generally much higher)
numerical scale than Table~\ref{tab:main_results}, and
\textbf{are not directly comparable to it}.
All comparisons in the following sections should be read only
relative to one another, within the same section.

\paragraph{Reading Precision and Recall together.}
Under this relaxed protocol, every one of the 101 queries has exactly
1\,000 relevant companies. This has a direct consequence for how
Precision@$k$ and Recall@$k$ relate to one another at any cut-off
$k \le 1\,000$: since both metrics are computed from the same count
(how many of the top-$k$ retrieved companies are relevant), just divided
by a different denominator,
\[
  \text{Recall@}k = \text{Precision@}k \times \frac{k}{1\,000},
\]
so the two metrics are exactly proportional to one another at every
cut-off, and neither is redundant on its own: Precision@$k$ answers
``of what the system returned, how much was correct'', while Recall@$k$
answers ``of everything that is correct, how much did the system
return'', and the two only coincide once $k = 1\,000$, where the
denominators become equal.
A practical consequence is that Recall@$k$ for small $k$ can look
misleadingly low in isolation. For example, a system that returns 10
companies of which 8 are genuinely relevant has an excellent
Precision@10 of 0.80, but the corresponding Recall@10 is only
$0.80 \times 10/1\,000 = 0.008$, not because the 10 companies returned
are poor, but because 10 items can never cover more than 1\% of a
1\,000-company relevant set. For this reason, the tables in this
chapter report Precision@$k$ at the shallow cut-offs
($k \in \{10, 50, 100, 300, 500\}$), where it is the more directly
interpretable number, and report both Precision@1\,000 and Recall@1\,000
together, since $k = 1\,000$ is the one cut-off at which Recall stops
being a rescaled copy of Precision and becomes a genuinely independent
measure: total coverage of the full relevant set. Recall@1\,000 is
treated as the headline metric throughout this thesis for exactly this
reason.

% ────────────────────────────────────────────────────────────
\section{Extended Baseline Comparison: Additional Embedding Models}
\label{sec:extended_baselines}
% ────────────────────────────────────────────────────────────

Five further dense embedding models were evaluated on the same
98\,716-company corpus, all-fields text representation, and 101-query
evaluation set described in Section~\ref{subsec:dataset}:

\begin{itemize}
    \item \textbf{BGE-M3} (\texttt{BAAI/bge-m3}): a multi-lingual,
    multi-functionality embedding model, 1\,024-dimensional embeddings.
    \item \textbf{GTE-large} (\texttt{thenlper/gte-large}): a
    general-purpose text embedding model, 1\,024 dimensions.
    \item \textbf{E5-Mistral-7B-Instruct}
    (\texttt{intfloat/e5-mistral-7b-instruct}): an instruction-tuned
    7-billion-parameter Mistral-based embedding model, 4\,096 dimensions.
    \item \textbf{Linq-Embed-Mistral}
    (\texttt{Linq-AI-Research/Linq-Embed-Mistral}): a 7B-parameter
    Mistral-based embedding model fine-tuned for retrieval,
    4\,096 dimensions.
    \item \textbf{SFR-Embedding-Mistral}
    (\texttt{Salesforce/SFR-Embedding-Mistral}): Salesforce's
    7B-parameter Mistral-based retrieval embedding model,
    4\,096 dimensions.
\end{itemize}

\begin{table}[ht]
\centering
\caption{Extended baseline Recall@$k$ under the relaxed (top-1\,000)
         pseudo-relevance protocol (Section~\ref{sec:protocol_note}).
         Bold indicates the best value in each column.}
\label{tab:extended_baselines}
\begin{tabular}{lccccc}
\toprule
& \multicolumn{4}{c}{\textbf{Precision@$k$}} & \textbf{Recall@1000} \\
\cmidrule(lr){2-5}
\textbf{Method} & \textbf{10} & \textbf{50} & \textbf{100} & \textbf{1000} & \\
\midrule
BGE-M3        & 0.927 & 0.916 & 0.905 & 0.704          & 0.704 \\
GTE-large     & 0.946 & 0.936 & 0.925 & \textbf{0.768} & \textbf{0.768} \\
E5-Mistral-7B & 0.629 & 0.572 & 0.523 & 0.274          & 0.274 \\
Linq-Mistral  & 0.934 & 0.921 & 0.915 & 0.749          & 0.749 \\
SFR-Mistral   & 0.933 & 0.922 & 0.915 & 0.746          & 0.746 \\
\bottomrule
\end{tabular}
\end{table}

\paragraph{Findings.}
GTE-large is the strongest of the five additional models
(Recall@1000~=~0.768, NDCG@10~=~0.946), ahead of Linq-Mistral
(0.749 / 0.936), SFR-Mistral (0.746 / 0.938), and BGE-M3
(0.704 / 0.935). These four sit close together and clearly ahead of
E5-Mistral-7B-Instruct, which substantially underperforms
(Recall@1000~=~0.274, NDCG@10~=~0.631) despite sharing the same base
architecture and parameter count as Linq-Mistral and SFR-Mistral.
This indicates that fine-tuning data and strategy, not model size, are
the dominant factor for this retrieval task. Notably, the two much
smaller 1\,024-dimensional models (BGE-M3, GTE-large) are competitive
with or better than the 7B-parameter models, at a fraction of the
encoding cost.

\paragraph{Encoding cost.}
The three 7B-parameter models required substantially more compute to
encode the corpus: E5-Mistral took 2\,011 seconds (33.5 min),
Linq-Mistral 2\,039 seconds (34.0 min), and SFR-Mistral 2\,032 seconds
(33.9 min), all on an A100 GPU, and all requiring checkpoint/resume
logic to fit within the cluster's 30-minute job wall-time limit.
This is roughly five times the encoding time of the original
335M-parameter BGE-large baseline (7.0 minutes, Section~\ref{subsec:bge}),
for retrieval quality that is not better and, in E5-Mistral's case,
substantially worse.

\paragraph{Incomplete and abandoned models.}
Two further large embedding models were attempted but did not produce
final results. \textbf{BGE-Multilingual-Gemma2}
(\texttt{BAAI/bge-multilingual-gemma2}, 9B parameters) encoding is
still in progress at the time of writing: model loading alone takes
15--27 minutes on this cluster's shared storage (highly variable due
to filesystem I/O contention), leaving little of the 30-minute job
window for actual encoding, so the job must be resubmitted repeatedly
to make incremental progress via checkpointing. \textbf{NV-Embed-v2}
(\texttt{nvidia/NV-Embed-v2}) was abandoned after a persistent
\texttt{AttributeError} caused by an incompatibility between the
model's custom code and the installed \texttt{transformers} library
version, which could not be resolved without a library downgrade that
risked breaking the other baselines' environment.

% ────────────────────────────────────────────────────────────
\section{BGE Instruction-Prefix Correction}
\label{sec:bge_prefix}
% ────────────────────────────────────────────────────────────

The original BGE-large baseline (Section~\ref{subsec:bge}) encoded
queries without an instruction prefix. BGE-large's model card
specifies that, for retrieval tasks, each \emph{query} should be
prefixed with the instruction \textit{``Represent this sentence for
searching relevant passages: ''} before encoding; the passage
(company) side requires no prefix, so the cached company embeddings
could be reused without re-encoding the corpus. This fix was retested
on both the all-fields and summary-only variants of the BGE baseline.

\begin{table}[ht]
\centering
\caption{Effect of the query instruction prefix on BGE-large, measured
         under the relaxed top-1\,000 protocol (Section~\ref{sec:protocol_note}).}
\label{tab:bge_prefix}
\begin{tabular}{lcccc}
\toprule
\textbf{Variant} & \textbf{NDCG@10} & \textbf{Recall@1000} & $\Delta$\textbf{NDCG@10} & $\Delta$\textbf{Recall@1000} \\
\midrule
All-fields, no prefix   & 0.942 & 0.726 & --      & --      \\
All-fields, with prefix & 0.937 & 0.727 & $-$0.005 & $+$0.001 \\
Summary-only, no prefix   & 0.937 & 0.700 & --      & --      \\
Summary-only, with prefix & 0.924 & 0.707 & $-$0.014 & $+$0.007 \\
\bottomrule
\end{tabular}
\end{table}

\paragraph{Finding.}
Contrary to expectation, applying the officially documented query
prefix produced only a marginal and mixed effect on this dataset:
Recall@1000 improved negligibly (+0.001 for all-fields, +0.007 for
summary-only), while NDCG@10 slightly \emph{decreased} in both variants.
Both effects are small enough to plausibly be within query-level noise.
This suggests that, for this company-retrieval task, BGE-large's query
prefix convention, tuned primarily for open-domain passage retrieval
benchmarks, does not transfer into a measurable improvement, and the
model's underlying limitation observed in Section~\ref{subsec:bge}
(underperforming MiniLM on Recall@1000) is not explained by the missing
prefix.

% ────────────────────────────────────────────────────────────
\section{Query Expansion: HyDE Pipeline}
\label{sec:hyde}
% ────────────────────────────────────────────────────────────

HyDE (Hypothetical Document Embeddings) expands each query by asking a
local LLM (\texttt{llama3.1:8b} via Ollama) to generate $N=10$
stochastic (temperature~$=0.8$) hypothetical company descriptions that
would ideally answer the query. The embeddings of these 10 descriptions
are averaged into a single query vector, which is then searched against
the MiniLM company index. A second stage reranks the top-50 HyDE
results with a cross-encoder (\texttt{BAAI/bge-reranker-v2-m3}).

\begin{table}[ht]
\centering
\caption{HyDE pipeline results (relaxed top-1\,000 protocol) compared
         against the plain MiniLM baseline re-evaluated under the same
         protocol.}
\label{tab:hyde}
\begin{tabular}{lcccc}
\toprule
\textbf{Method} & \textbf{NDCG@10} & \textbf{Recall@1000} & \textbf{Latency (ms)} \\
\midrule
MiniLM (no expansion)        & 0.975 & 0.741 & 16.9 \\
Stage 1: HyDE + MiniLM       & 0.829 & 0.596 & 9\,246.9 \\
Stage 2: HyDE + MiniLM + Reranker & 0.954 & 0.596 & 9\,488.3 \\
\bottomrule
\end{tabular}
\end{table}

\paragraph{Finding.}
Query expansion via HyDE \emph{hurts} retrieval quality relative to
directly encoding the original query: Stage 1 NDCG@10 drops from 0.975
to 0.829, and Recall@1000 drops from 0.741 to 0.596. Adding a
cross-encoder reranker (Stage 2) recovers most of the NDCG@10 loss
(0.829~$\to$~0.954) by re-ordering the top-50 candidates, but cannot
recover the recall loss, since reranking only reorders an already
-fixed candidate set. Latency is dominated overwhelmingly by the LLM
call itself (9.2 seconds), three orders of magnitude slower than
directly encoding the query with MiniLM (16.9~ms). For this dataset,
HyDE is not a worthwhile query-expansion strategy: it is far slower and
retrieves a worse candidate set than simply embedding the query
directly.

% ────────────────────────────────────────────────────────────
\section{Hybrid Retrieval: Reciprocal Rank Fusion}
\label{sec:rrf}
% ────────────────────────────────────────────────────────────

As an alternative to query expansion, BM25 and MiniLM rankings are
fused directly using Reciprocal Rank Fusion (RRF),
$\mathrm{RRF}(d) = \sum_{\text{method}} \frac{1}{k + \mathrm{rank}(d)}$
with $k=60$ \citep{cormack2009reciprocal}. A second stage reranks the
top-50 fused results with the same cross-encoder used for HyDE.

\begin{table}[ht]
\centering
\caption{Hybrid RRF results (relaxed top-1\,000 protocol).}
\label{tab:rrf}
\begin{tabular}{lcccc}
\toprule
\textbf{Method} & \textbf{NDCG@10} & \textbf{Recall@1000} & \textbf{Latency (ms)} \\
\midrule
MiniLM (no fusion)              & 0.975 & 0.741 & 16.9 \\
Stage 1: RRF (BM25 + MiniLM)    & 0.955 & 0.724 & 104.9 \\
Stage 2: RRF + Reranker          & 0.975 & 0.724 & 352.0 \\
\bottomrule
\end{tabular}
\end{table}

\paragraph{Finding.}
Unlike HyDE, RRF fusion of BM25 and MiniLM adds negligible latency
(the fusion step itself takes 2.2\,ms) and, after reranking, matches
the plain MiniLM baseline's NDCG@10 (0.975) while Recall@1000 remains
slightly below MiniLM alone (0.724 vs 0.741): BM25's keyword matches
do not add companies beyond what MiniLM already retrieves for this
query set, so fusion mostly re-confirms rather than extends MiniLM's
candidate pool. RRF is a much cheaper and more effective way to
incorporate a second retrieval signal than HyDE-style query expansion.

% ────────────────────────────────────────────────────────────
\section{LLM-as-Judge Relevance Validation}
\label{sec:llm_judge}
% ────────────────────────────────────────────────────────────

Because the pseudo-relevance labels (Section~\ref{subsec:eval_protocol})
assume production's own rankings are correct, an independent check was
run using an LLM to directly judge relevance. For 20 held-out queries
(from a pool of 101, split 81/20 train/held-out), the top-20 production
results per query were judged individually on a 0/1/2 relevance scale
by an LLM given the query and the company summary, yielding 400
graded judgements in total.

\begin{table}[ht]
\centering
\caption{NDCG under pseudo-relevance labels vs.\ LLM-judged relevance
         labels, on the 20 held-out queries.}
\label{tab:llm_judge}
\begin{tabular}{lccc}
\toprule
\textbf{$k$} & \textbf{Pseudo-relevance NDCG} & \textbf{LLM-judged NDCG} & $\Delta$ \\
\midrule
10   & 0.964 & 0.963 & $-$0.001 \\
50   & 0.947 & 0.982 & $+$0.035 \\
100  & 0.939 & 0.982 & $+$0.042 \\
500  & 0.868 & 0.982 & $+$0.113 \\
1000 & 0.771 & 0.982 & $+$0.210 \\
\bottomrule
\end{tabular}
\end{table}

\paragraph{Finding.}
At $k=10$, pseudo-relevance and LLM-judged NDCG agree closely
(0.964 vs 0.963), supporting the use of production's own top results
as a reasonable relevance proxy at shallow depths. The two diverge
sharply as $k$ grows: LLM-judged NDCG stays essentially flat at 0.982
out to $k=1000$, while pseudo-relevance NDCG falls to 0.771. This is
consistent with the pseudo-relevance labels being a hard top-100 cutoff
(Section~\ref{subsec:eval_protocol}): once retrieval looks past that
cutoff, pseudo-relevance necessarily penalises the system even for
companies an LLM judge considers genuinely relevant. This supports the
interpretation, already flagged in Section~\ref{subsec:eval_protocol},
that pseudo-relevance NDCG at large $k$ understates true system quality
rather than reflecting a real quality gap.

% ────────────────────────────────────────────────────────────
\section{Production API Evaluation and MiniLM Rerank}
\label{sec:api_eval}
% ────────────────────────────────────────────────────────────

Two rounds of evaluation compared the local pipeline against Istari's
live production search API, using a separate ground truth (the
production system's own graded semantic-search results at a similarity
threshold of 0.75) rather than the pseudo-relevance labels used
elsewhere in this chapter.

\paragraph{v1: keyword API.}
The first round compared Istari's v1 (BM25-based) production API
against local MiniLM, both queried against the same 101 queries.
Istari v1 BM25 NDCG@10 was 0.101 and Recall@1000 was 0.109 (network
latency 2\,828\,ms/query), against local MiniLM's NDCG@10 of 0.935 and
Recall@1000 of 0.766 (local GPU latency 26.6\,ms/query), confirming
that the production keyword API is a much weaker retrieval signal than
even the simplest local dense baseline.

\paragraph{v2: semantic API.}
A second round added Istari's newer v2 semantic search API (live,
paginated over the full 20-million-company GOI, quota-aware). Istari
v2 NDCG@10 was 0.985 and Recall@1000 (at the achievable depth of 816
results) was 0.650, clearly ahead of local MiniLM (NDCG@10 0.935,
Recall@816 0.709) at shallow cutoffs, though local MiniLM's Recall
catches up and slightly exceeds v2's at $k=816$. Network latency for
v2 (1\,966\,ms/query for a full fetch) remains far higher than local
MiniLM (26.4\,ms/query).

\paragraph{Reranking v2's own candidates with MiniLM.}
A follow-up experiment took v2's own retrieved candidate pool and
re-ordered it locally using MiniLM, to test whether local reranking
could improve on v2's native ranking without changing which companies
are retrieved. This made results \emph{worse}: NDCG@10 fell from 0.985
(v2 native ranking) to 0.855 (MiniLM-reranked). Istari's v2 semantic
ranking is not simply a MiniLM-style dense embedding search under the
hood: it outperforms MiniLM's own ranking of the same candidates.

% ────────────────────────────────────────────────────────────
\section{ANN Index Scaling to 2 Million Vectors}
\label{sec:ann_scaling}
% ────────────────────────────────────────────────────────────

To test whether the exact search used throughout this chapter
(\texttt{IndexFlatL2}/\texttt{IndexFlatIP}) remains practical at
production scale, the 98\,716 real MiniLM company embeddings were
padded with 1\,901\,284 synthetic filler vectors (real embeddings
perturbed with Gaussian noise at 15\% of the average vector norm),
giving a 2-million-vector corpus. Exact search (\texttt{FlatL2}) was
compared against approximate indexes: IVF (\texttt{nlist}=4096, with
\texttt{nprobe} $\in \{1, 8, 32, 128, 4096\}$) and HNSW ($M=32$, with
\texttt{efSearch} $\in \{16, 64, 256\}$).

\begin{table}[ht]
\centering
\caption{Exact vs.\ approximate search at 2M-vector scale (relaxed
         top-1\,000 protocol). Only representative configurations shown.}
\label{tab:ann_scaling}
\begin{tabular}{lccccc}
\toprule
\textbf{Index} & \textbf{NDCG@10} & \textbf{Recall@1000} & \textbf{ms/query} & \textbf{Speedup} & \textbf{Index size} \\
\midrule
Exact (FlatL2)        & 0.975 & 0.741 & 175.5 & 1.0$\times$ & 3\,072\,MB \\
IVF (nprobe=8)         & 0.974 & 0.743 & 26.4  & 6.6$\times$ & 3\,094\,MB \\
HNSW (efSearch=256)    & 0.975 & 0.727 & 26.9  & 6.5$\times$ & 3\,616\,MB \\
\bottomrule
\end{tabular}
\end{table}

\paragraph{Finding.}
Both IVF (\texttt{nprobe}=8) and HNSW (\texttt{efSearch}=256) recover
essentially all of exact search's NDCG@10, and IVF slightly exceeds
exact search's Recall@1000 (0.743 vs 0.741, within noise), while
running 6.5--6.6$\times$ faster at 2-million-vector scale, at the cost
of modest additional index memory (0.7\% and 17.7\% larger,
respectively). This confirms that approximate search is a safe and
substantial efficiency gain for scaling the retrieval system toward
the full 20-million-company GOI, without a meaningful quality trade-off
at this scale.

% ────────────────────────────────────────────────────────────
\section{Recall@1000 Error Analysis}
\label{sec:error_analysis}
% ────────────────────────────────────────────────────────────

To understand why exact MiniLM search caps at Recall@1000~=~0.741
(Section~\ref{sec:ann_scaling}), every production-relevant company that
MiniLM failed to retrieve within the top-1\,000 was extracted and
searched for at much greater depth ($k=10\,000$), yielding
26\,185 missed (query, company) pairs across all 101 queries.

\begin{table}[ht]
\centering
\caption{Classification of the 26\,185 Recall@1000 misses by how far
         beyond rank 1\,000 the true company was actually ranked.}
\label{tab:error_analysis}
\begin{tabular}{lrr}
\toprule
\textbf{Bucket} & \textbf{Count} & \textbf{Share (\%)} \\
\midrule
Near miss (rank 1\,000--2\,000)   & 16\,157 & 61.7 \\
Mid miss (rank 2\,000--5\,000)    &  8\,061 & 30.8 \\
Far miss (rank 5\,000--10\,000)   &  1\,375 &  5.3 \\
Not found (beyond rank 10\,000)   &    592 &  2.3 \\
\bottomrule
\end{tabular}
\end{table}

\paragraph{Finding.}
The large majority of misses (61.7\%) are near misses, just outside the
top-1\,000 cutoff, suggesting most of the recall gap is a matter of
ranking granularity rather than a hard semantic failure. Only 2.3\% of
misses are not found at all within $k=10\,000$, indicating genuine
embedding failures are rare. Separately, 3.8\% of missed companies
(996 of 26\,185) were flagged as boilerplate content, and one extreme
case of description duplication was found: the placeholder text
``Bluehost.com'' appears as the company description for 709 distinct
domains, meaning these domains are indistinguishable to any
text-embedding model and can never be individually retrieved correctly.
This is a concrete data-quality artifact rather than a retrieval-model
limitation, and sets a hard ceiling on achievable recall for any
embedding-based method operating on the \texttt{summary} field alone.

% ────────────────────────────────────────────────────────────
\section{Oracle Union Recall Ceiling}
\label{sec:oracle_union}
% ────────────────────────────────────────────────────────────

Before investing in a trained reranker or a fine-tuning pipeline, it is
useful to know the best-case recall achievable by simply combining the
candidate pools already produced by every method run so far --- at zero
additional encoding cost, since every channel's top-1\,000 results per
query are already cached on disk. This experiment computes, for every
combination of the six available channels (MiniLM, BM25, BGE-all-fields,
BGE-summary, OpenAI-large, Nomic), the recall of the \emph{union} of
their top-1\,000 candidate sets against the top-1\,000 ground truth.

\paragraph{Findings.}
The best single channel (MiniLM alone) reaches Recall@1000~=~0.741.
The full union of all six channels reaches 0.936 (average combined pool
size 1\,949 candidates/query), a gain of +0.196 over MiniLM alone. The
best-performing \emph{pair} is MiniLM~+~OpenAI-large, at 0.856 (average
pool size 1\,300). A pool-size sweep on this best pair shows the
ceiling is highly sensitive to how many candidates are taken per
channel: at 100 candidates/channel the union only reaches 0.150 recall,
rising to 0.581 at 500, and 0.856 at the full 1\,000.

\paragraph{Interpretation.}
Because the six-channel union clears 0.936, well above any individual
channel, the different retrieval methods are largely making
\emph{complementary}, not correlated, errors: a reranker operating over
a modestly sized union pool (rather than a single channel) is likely
sufficient to approach this ceiling, and a full fine-tuning pipeline is
not obviously necessary to close the remaining gap. The residual
0.064 recall lost even under the full six-channel union is consistent
with the hard data-quality ceiling identified in
Section~\ref{sec:error_analysis} (duplicated/boilerplate company
descriptions that no embedding-based channel can disambiguate).

% ────────────────────────────────────────────────────────────
\section{SimLM: A Purpose-Built Retrieval Pre-training Baseline}
\label{sec:simlm}
% ────────────────────────────────────────────────────────────

All dense baselines evaluated so far are general-purpose sentence-embedding
models. \textbf{SimLM} \citep{wang2022simlm} takes a different approach:
it is a BERT-base (110\,M parameter) encoder pre-trained specifically for
passage retrieval using a representation-bottleneck objective, before any
task-specific fine-tuning. It is included to test whether a model
purpose-built for retrieval, but far smaller than the 7B-parameter models
in Section~\ref{sec:extended_baselines}, can match or exceed them.

\paragraph{Implementation.}
The public checkpoint \texttt{intfloat/simlm-base-msmarco-finetuned} is
used, following its documented convention rather than the all-fields
concatenation used elsewhere in this thesis: the company name is encoded
as a \emph{title}, and the remaining fields (country, state, city, type,
size, industry, summary, keywords) as the paired \emph{body} text, using
CLS-token pooling (not mean pooling) with L2 normalisation. Queries are
encoded without any prefix, following the same model card.

\begin{table}[ht]
\centering
\caption{SimLM Recall@1000 (relaxed top-1\,000 protocol) alongside the
         strongest baselines evaluated so far, for reference.}
\label{tab:simlm}
\begin{tabular}{lccccc}
\toprule
& \multicolumn{4}{c}{\textbf{Precision@$k$}} & \textbf{Recall@1000} \\
\cmidrule(lr){2-5}
\textbf{Method} & \textbf{10} & \textbf{50} & \textbf{100} & \textbf{1000} & \\
\midrule
BM25 (Section~\ref{sec:baseline_sparse})         & 0.838 & 0.748 & 0.723 & 0.635 & 0.635 \\
SimLM                                            & 0.920 & 0.881 & 0.861 & 0.609 & 0.609 \\
MiniLM (Section~\ref{subsec:minilm})             & 0.972 & 0.950 & 0.938 & 0.741 & 0.741 \\
GTE-large (Section~\ref{sec:extended_baselines}) & 0.946 & 0.936 & 0.925 & \textbf{0.768} & \textbf{0.768} \\
\bottomrule
\end{tabular}
\end{table}

\paragraph{Finding.}
SimLM underperforms every dense baseline evaluated in this thesis, and its
Recall@1000 (0.609) falls even below the lexical BM25 baseline (0.635),
despite SimLM being purpose-pretrained for passage retrieval. (NDCG@10 tells
a less clear-cut story here: SimLM's 0.923 sits above BM25's 0.854, since
NDCG rewards ranking quality among whichever companies are retrieved, not
coverage of the full relevant set.) A likely
explanation is a domain mismatch between SimLM's pre-training convention
and this task: SimLM was pre-trained and fine-tuned on MS~MARCO, a
corpus of short, single-topic passages, using a title/body sentence-pair
format that assumes the title is a concise topical label. Company names
are not topical labels in the same sense: ``Sweet Software'' or
``Springcoast Capital Partners'' convey little about what the paired body
text describes, so the sentence-pair convention that helps SimLM on
MS~MARCO-style data may provide a weaker signal here than simply
concatenating all fields into one string, as the other baselines do.
This result is a useful negative finding: retrieval-specific pre-training
does not automatically transfer across domains, and general-purpose
sentence embeddings (MiniLM, GTE-large) remain the stronger choice for
this task.

% ────────────────────────────────────────────────────────────
\section{ColBERT: Late-Interaction Retrieval from Scratch}
\label{sec:colbert}
% ────────────────────────────────────────────────────────────

Every method evaluated so far represents each company and each query as a
single dense vector. \textbf{ColBERT} \citep{khattab2020colbert,
santhanam2022colbertv2} instead represents both as \emph{sets of
per-token} embeddings and scores a (query, company) pair using
\emph{MaxSim}: for every query token, the maximum similarity to any
company token is taken, and these per-token maxima are summed. This is
architecturally richer than single-vector dense retrieval, since no
single fixed-size vector has to compress the entire meaning of a company
description, and it is included to test whether that added expressiveness
translates into higher recall on this task.

\paragraph{Implementation.}
The public checkpoint \texttt{colbert-ir/colbertv2.0} is used, encoding
each company summary into a sequence of 128-dimensional per-token
embeddings via the model's linear projection layer. Due to indexing and
library constraints on the cluster used for this thesis, the official
PLAID-indexed ColBERT pipeline was not used; instead, MaxSim scoring is
implemented from scratch as a brute-force computation over all
98\,716 companies for every query. This is computationally more expensive
than an indexed implementation would be, but is exact rather than
approximate, and is tractable at this corpus size.

\begin{table}[ht]
\centering
\caption{ColBERT (brute-force MaxSim) Recall@1000 (relaxed top-1\,000
         protocol) alongside the strongest single-vector baselines.}
\label{tab:colbert}
\begin{tabular}{lccccc}
\toprule
& \multicolumn{4}{c}{\textbf{Precision@$k$}} & \textbf{Recall@1000} \\
\cmidrule(lr){2-5}
\textbf{Method} & \textbf{10} & \textbf{50} & \textbf{100} & \textbf{1000} & \\
\midrule
MiniLM (Section~\ref{subsec:minilm})             & 0.972 & 0.950 & 0.938 & 0.741 & 0.741 \\
GTE-large (Section~\ref{sec:extended_baselines}) & 0.946 & 0.936 & 0.925 & 0.768 & 0.768 \\
ColBERT (MaxSim)                                 & 0.936 & 0.931 & 0.925 & \textbf{0.773} & \textbf{0.773} \\
\bottomrule
\end{tabular}
\end{table}

\paragraph{Finding.}
ColBERT's Recall@1000 (0.773) is the highest of any single retrieval
channel evaluated in this thesis, narrowly ahead of GTE-large (0.768) and
clearly ahead of MiniLM (0.741). Its NDCG@10 (0.937), however, is lower
than both MiniLM (0.975) and GTE-large (0.946), meaning ColBERT retrieves
a broader set of relevant companies overall but ranks the very top results
somewhat less precisely. Token-level interaction therefore does deliver a
genuine, if modest, recall improvement over the best single-vector models,
though not one large enough on its own to close the gap to the hybrid
rerank pipelines described in Section~\ref{sec:hybrid_fullpool}.

% ────────────────────────────────────────────────────────────
\section{Hybrid Full-Pool Reranking}
\label{sec:hybrid_fullpool}
% ────────────────────────────────────────────────────────────

Section~\ref{sec:rrf} showed that fusing BM25 and MiniLM with Reciprocal
Rank Fusion, followed by reranking only the top-50 fused candidates,
essentially reproduces MiniLM's own quality (NDCG@10~=~0.975) without
improving Recall@1000. Section~\ref{sec:oracle_union} then showed that the
\emph{union} of several channels' top-1\,000 candidate pools reaches a
much higher oracle recall ceiling than any single channel. Together, these
two findings motivate a different hybrid design: rerank the
\emph{entire} union pool of several dense channels with a cross-encoder,
rather than fusing rankings and reranking only a shallow, fixed-size
slice of the result.

\paragraph{Implementation.}
For each query, the top-1\,000 results from a set of dense channels are
combined by simple set union (no rank fusion, no truncation), giving an
unranked candidate pool that is typically 1\,300--1\,800 companies. A
cross-encoder reranker, \texttt{BAAI/bge-reranker-v2-m3}
\citep{baaireranker2024}, then scores every (query, candidate) pair in
the full pool directly, and the top-1\,000 by reranker score are kept as
the final result. Four variants are compared, adding one channel at a
time:

\begin{itemize}
    \item \textbf{2-way}: MiniLM~+~Linq-Embed-Mistral.
    \item \textbf{3-way}: 2-way~+~GTE-large.
    \item \textbf{4-way}: 3-way~+~OpenAI-large.
\end{itemize}

\begin{table}[ht]
\centering
\caption{Full-pool reranking (relaxed top-1\,000 protocol), compared
         against the shallow RRF~+~top-50-rerank baseline from
         Section~\ref{sec:rrf}. Bold indicates the best value.}
\label{tab:hybrid_fullpool}
\small
\setlength{\tabcolsep}{2.5pt}
\begin{tabular}{lccccc}
\toprule
\textbf{Method} & \textbf{P@10} & \textbf{P@100} & \textbf{P@1000} & \textbf{R@1000} & \textbf{Pool} \\
\midrule
RRF + top-50 (\S\ref{sec:rrf})   & 0.980 & 0.940 & 0.782          & 0.782          & 50 \\
2-way (MiniLM+Linq)              & 0.955 & 0.950 & 0.802          & 0.802          & 1\,300 \\
3-way (+GTE-large)               & 0.955 & 0.950 & \textbf{0.804} & \textbf{0.804} & 1\,424 \\
4-way (+OpenAI-large)            & 0.955 & 0.950 & 0.802          & 0.802          & 1\,600 \\
\bottomrule
\end{tabular}
\normalsize
\end{table}

\paragraph{Finding.}
Reranking the full union pool instead of a fixed top-50 slice raises
Recall@1000 substantially (0.782~$\to$~0.802, 2-way), at the cost of a
small drop in NDCG@10 (0.979~$\to$~0.955): the reranker now has many
more genuinely difficult candidates to rank, so its ordering of the very
top results is slightly less precise, even though it recovers more of the
true relevant set overall. Adding a third channel (GTE-large) improves
Recall@1000 further, to 0.804, the best result obtained by any hybrid
pipeline built from off-the-shelf pretrained components in this thesis.
Adding a fourth channel (OpenAI-large) does \emph{not} improve on this --
Recall@1000 is unchanged at 0.802, within noise of the 3-way result.

\paragraph{Interpretation: channel saturation.}
This flat 3-way~$\to$~4-way result is an important turning point in this
thesis's experimental narrative. Section~\ref{sec:oracle_union} showed
that combining \emph{more} candidate-generating channels reliably raises
the oracle recall ceiling; here, adding a fourth channel to an already
capable 3-way pool does not translate into a better final result. This
suggests that candidate coverage is no longer the limiting factor for this
pipeline: the union pool already contains the relevant companies needed
to reach a substantially higher recall than 0.804 (recall
Section~\ref{sec:oracle_union}'s six-channel oracle ceiling of 0.936) --
and the reranker itself is now the bottleneck preventing the pipeline from
closing that gap. This motivates two follow-up questions, addressed in the
next two sections: is the specific reranker \emph{model} the limiting
factor (Section~\ref{sec:reranker_ablation}), or is a fundamentally
different fusion approach needed (Section~\ref{sec:learned_fusion})?

% ────────────────────────────────────────────────────────────
\section{Reranker Ablation: Isolating the Bottleneck}
\label{sec:reranker_ablation}
% ────────────────────────────────────────────────────────────

Given the channel-saturation finding in Section~\ref{sec:hybrid_fullpool},
the natural next question is whether the specific reranker model used
throughout this chapter, \texttt{BAAI/bge-reranker-v2-m3}, is itself
limiting recall, or whether reranker choice matters little once the
candidate pool is fixed.

\paragraph{Implementation.}
\texttt{BAAI/bge-reranker-v2-m3} (built on an XLM-RoBERTa backbone) is
replaced with \texttt{mixedbread-ai/mxbai-rerank-large-v1}
\citep{mixedbread2024mxbai} (built on a DeBERTa-v2 backbone, an
architecturally distinct, comparably-sized cross-encoder), applied to the
\emph{exact same} 3-way union pool from Section~\ref{sec:hybrid_fullpool}.
Only the reranker model changes; the candidate pool, corpus, and queries
are identical.

\begin{table}[ht]
\centering
\caption{Reranker ablation Recall@1000: identical 3-way candidate pool,
         two different reranker models (relaxed top-1\,000 protocol).}
\label{tab:reranker_ablation}
\small
\setlength{\tabcolsep}{2.5pt}
\begin{tabular}{lcccc}
\toprule
\textbf{Reranker} & \textbf{P@10} & \textbf{P@100} & \textbf{P@1000} & \textbf{R@1000} \\
\midrule
BAAI/bge-reranker-v2-m3 (\S\ref{sec:hybrid_fullpool}) & 0.955 & 0.950 & \textbf{0.8042} & \textbf{0.8042} \\
mxbai-rerank-large-v1                                 & 0.968 & 0.940 & 0.8035          & 0.8035 \\
\bottomrule
\end{tabular}
\normalsize
\end{table}

\paragraph{Finding.}
Recall@1000 is essentially unchanged between the two rerankers (0.8042
vs.\ 0.8035, a difference of 0.0007, well within query-level noise).
NDCG@10 does improve modestly with the alternative reranker (0.955~$\to$~0.971),
indicating that \texttt{mxbai-rerank-large-v1} orders the very top
results somewhat better, without changing which companies land in the
top-1\,000 overall. The conclusion is that reranker \emph{model choice} is
not the bottleneck identified in Section~\ref{sec:hybrid_fullpool}: the
remaining gap to the oracle ceiling is not a matter of picking a better
off-the-shelf cross-encoder. This points towards a more fundamental
limitation: either genuinely hard, vocabulary-mismatched cases that no
cross-encoder scoring the \texttt{summary} field alone can resolve, or a
fusion mechanism (simple set union plus independent reranking) that is
itself insufficiently expressive, motivating the learned approach in
Section~\ref{sec:learned_fusion}.

% ────────────────────────────────────────────────────────────
\section{Learned Fusion Ranker}
\label{sec:learned_fusion}
% ────────────────────────────────────────────────────────────

Every hybrid pipeline so far combines \emph{frozen}, independently
pretrained components: candidate lists are pooled by simple union, and a
pretrained cross-encoder, never fit to this task's own queries or
labels, reranks the result. This section instead treats fusion as a
supervised learning-to-rank problem \citep{liu2009learning}: a model is
trained directly on this thesis's own 101 queries and their pseudo-relevance
labels to learn how to combine multiple retrieval signals, rather than
relying on a fixed, hand-designed combination rule.

\paragraph{Implementation.}
For every (query, candidate) pair in the same 3-way union pool used in
Section~\ref{sec:hybrid_fullpool} (with BM25 added as a fourth,
lexical, signal), ten features are computed: the raw similarity score
and inverse rank ($1/\text{rank}$) from each of MiniLM, Linq-Embed-Mistral,
GTE-large, and BM25; the \texttt{BAAI/bge-reranker-v2-m3} score from
Section~\ref{sec:hybrid_fullpool} (included as one input feature among
ten, not as the final decision rule); and a count of how many channels
retrieved the candidate at all. Two models are trained on these features
using \texttt{scikit-learn} \citep{pedregosa2011scikit}: a logistic
regression, and a histogram-based gradient-boosted tree ensemble.
Crucially, both are trained and evaluated using 5-fold cross-validation
grouped by \texttt{query\_id}: for each fold, the model is fit on
candidates from 80\% of queries and scores held-out candidates from the
remaining 20\%, so that no query is ever scored by a model that has seen
its own label during training. Final metrics are computed by
concatenating the held-out predictions from all five folds, giving a
complete, leakage-free evaluation over all 101 queries.

\begin{table}[ht]
\centering
\caption{Learned fusion ranker Recall@1000 (held-out, cross-validated
         predictions; relaxed top-1\,000 protocol), compared against the
         best pretrained-component pipeline from Section~\ref{sec:hybrid_fullpool}.
         Bold indicates the best value.}
\label{tab:learned_fusion}
\small
\setlength{\tabcolsep}{2.5pt}
\begin{tabular}{lcccc}
\toprule
\textbf{Method} & \textbf{P@10} & \textbf{P@100} & \textbf{P@1000} & \textbf{R@1000} \\
\midrule
3-way + BGE reranker (\S\ref{sec:hybrid_fullpool}) & 0.955 & 0.950 & 0.804          & 0.804          \\
Learned fusion, logistic regression                  & 0.968 & 0.953 & 0.821          & 0.821          \\
Learned fusion, gradient-boosted trees                & \textbf{0.977} & \textbf{0.955} & \textbf{0.823} & \textbf{0.823} \\
\bottomrule
\end{tabular}
\normalsize
\end{table}

\paragraph{Finding.}
The gradient-boosted learned fusion ranker reaches Recall@1000~=~0.823,
the best result obtained anywhere in this thesis, improving on the best
pretrained-component pipeline (0.804) by $+0.019$. For scale, this gain is
roughly 27 times larger than the 0.0007 movement observed when merely
swapping reranker models in Section~\ref{sec:reranker_ablation}; this is
a substantive improvement attributable to training on the task's own
data, not noise. NDCG@10 also improves (0.955~$\to$~0.979).

\paragraph{Interpretation: what the model learned.}
The trained logistic regression's coefficients offer a direct, interpretable
view into which signals matter most for this task: MiniLM's \emph{inverse
rank} has by far the largest coefficient magnitude (+5.01), exceeding
every raw similarity score \emph{and} the pretrained reranker's own score
(+0.16). In other words, where a candidate is ranked by MiniLM is a more
informative signal for final relevance than the pretrained cross-encoder's
score: a concrete, non-obvious finding that would not be visible from
simply combining components by hand. This supports the conclusion drawn
in Section~\ref{sec:reranker_ablation}: the pretrained reranker was not
extracting the maximum available signal from the candidate pool, and a
lightweight model trained on this task's own data closes a meaningful part
of the remaining gap to the oracle ceiling from Section~\ref{sec:oracle_union}.

% here the original content of my thesis ends

% This is the template here:
\chapter{Introduction}
\label{ch:intro}

The structure of this example thesis is \emph{non-binding}, but is highly
topic-dependent and should be discussed with your advisor. You probably do not
want to go below the level of subsections, though.

To cite, do it in one of the following ways:
\begin{itemize}
\item A single paper in text: As shown by \citet{doe2024proof}, X holds.
\item Same with additional information: As shown by \citet[p. 20]{doe2024proof},
  X holds.
\item A single paper in parenthesis: X holds \cite{doe2024proof}.
\item Multiple papers in parenthesis: X holds
  \cite{brown2022techreport,smith2023conference,lee2023journal,doe2024proof}.
\item Use footnotes to reference a website, such as the DWS thesis
  guidelines.\footnote{\url{https://www.uni-mannheim.de/dws/teaching/thesis-guidelines/}}
  Note that footnote marks are set after punctuation in English text. Only
  reference a website when needed---e.g., the website of a library or a blog
  post---and always cite corresponding scientific publications (if any) in
  addition.
\end{itemize}
Format the bibliography consistently, e.g., as in the the example bibliography
in this template. Only cite archival versions of papers (such as those published
on \href{https://arxiv.org}{arXiv}) when there are no corresponding official
publications (such as in conference proceedings); if there are such
publications, then cite those instead.


Some examples of references are:
\begin{enumerate}
\item Table~\ref{tab:sample} is an example table.
\item Section~\ref{sec:preliminaries} is a later section of this thesis.
\item Algorithm~\ref{alg:proofX} is an example algorithm.
\item And finally, Equation~\eqref{eq:s} is an example equation.
\end{enumerate}

Here is the equation referred to above:
\begin{equation}
  \label{eq:s}
  \mathcal{S}=\{\,a\cdot \operatorname{relu}(\mathbf{W^\top}\mathbf{x}) :
  a\in\mathbb{R} \,\}
\end{equation}

\textbf{Important:} Fill out and sign the last page when you submit/deliver the
final version of your thesis. Otherwise, your work cannot be accepted for legal
reasons.

The purpose of the introduction as summarized by~\citet{zobel2004}:

\blockcquote{zobel2004}{%
  An introduction can be regarded as an expanded version of the abstract. It
  should describe the paper's topic, the problem being studied, references to
  key papers, the approach to the solution, the scope and limitations of the
  solution, and the outcomes. There needs to be enough detail to allow readers
  to decide whether or not they need to read further. It should include
  motivation: the introduction should explain why the problem is interesting,
  what the relevant scientific issues are, and why the solution is a good one.

  That is, the introduction should show that the paper is worth reading and it
  should allow the reader to understand your perspective, so that the reader and
  you can proceed on a basis of common understanding.

  Many introductions follow a five-element organization:
  \begin{enumerate}
  \item A general statement introducing the broad research area of the
    particular topic being investigated.
  \item An explanation of the specific problem (difficulty, obstacle, challenge)
    to be solved.
  \item A brief review of existing or standard solutions to this problem and
    their limitations.
  \item An outline of the proposed new solution.
  \item A summary of how the solution was evaluated and what the outcomes of the
    evaluation were.
  \end{enumerate}

  An interesting exercise is to read other papers, analyze their introductions
  to see if they have this form, and then decide whether they are effective. The
  introduction can discuss the importance or ramifications of the conclusions
  but should omit supporting evidence, which the interested reader can find in
  the body of the paper. Relevant literature can be cited in the introduction,
  but unnecessary jargon, complex mathematics, and in-depth discussion of the
  literature belong elsewhere.

  A paper isn't a story in which results are kept secret until a surprise
  ending. The introduction should clearly tell the reader what in the paper is
  new and what the outcomes are. There may still be a little suspense: revealing
  what the results are does not necessarily reveal how they were achieved. If,
  however, the existence of results is concealed until later on, the reader
  might assume there are no results and discard the paper as worthless.}

\chapter{Literature Review}
\label{ch:related_work}

\citet{zobel2004} writes:

\blockcquote{zobel2004}{%
  Few results or experiments are entirely new. Most often they are extensions of
  or corrections to previous research-that is, most results are an incremental
  addition to existing knowledge. A literature review, or survey, is used to
  compare the new results to similar previously published results, to describe
  existing knowledge, and to explain how it is extended by the new results. A
  survey can also help a reader who is not expert in the field to understand the
  paper and may point to standard references such as texts or survey articles.

  In an ideal paper, the literature review is as interesting and thorough as the
  description of the paper's contribution. There is great value for the reader
  in a precise analysis of previous work that explains, for example, how
  existing methods differ from one another and what their respective strengths
  and weaknesses are. Such a review also creates a specific expectation of what
  the contribution of the paper should be-it shapes what the readers expect of
  your work, and thus shapes how they will respond to your ideas.

  The literature review can be early in a paper, to describe the context of the
  work, and might in that case be part of the introduction; or the literature
  review can follow or be part of the main body, at which point a detailed
  comparison between the old and the new can be made. If the literature review
  is late in a paper, it is easier to present the surveyed results in a
  consistent terminology, even when the cited papers have differing nomenclature
  and notation. In many papers the literature review material is not gathered
  into a single section, but is discussed where it is used-background material
  in the introduction, analysis of other researchers' work as new results are
  introduced, and so on. This approach can help you to write the paper as a
  flowing narrative.

  An issue that is difficult in some research is the relationship between new
  scientific results and proprietary commercial technology. It often is the case
  that scientists investigate problems that appear to be solved or addressed in
  commercial products. For example, there is ongoing academic research into
  methods for information retrieval despite the success of the search engines
  deployed on the web. From the perspective of high research principle, the
  existence of a commercial product is irrelevant: the ideas are not in the
  public domain, it is not known how the problems were solved in the product,
  and the researcher's contribution is valid. However, it may well be reckless
  to ignore the product; it should be cited and discussed, while noting, for
  example, that the methods and effectiveness of the commercial solution are
  unknown. }

An example structure for the literature review is given below; as before, this
structure is non-binding and should be changed as appropriate.

\section{Preliminaries}
\label{sec:preliminaries}

\section{Related Work on A}
\label{sec:related_work_A}

\section{Related Work on B}

\section{Related Work on C}

\section{Summary}

\chapter{Body Chapter 1}
\citet{zobel2004} writes:

\blockcquote{zobel2004}{%
  The body of a paper should present the results. The presentation should
  provide necessary background and terminology, explain the chain of reasoning
  that leads to the conclusions, provide the details of central proofs,
  summarize any experimental outcomes, and state in detail the conclusions
  outlined in the introduction. Descriptions of experiments should permit
  reproduction and verification, as discussed in Chapter 11. There should also
  be careful definitions of the hypothesis and major concepts, even those
  described informally in the introduction. The structure should be evident in
  the section headings. Since the body can be long, narrative flow and a clear
  logical structure are essential.

  The body should be reasonably independent of other papers. If, to understand
  your paper, the reader must find specialized literature such as your earlier
  papers or an obscure paper by your advisor, then its audience will be limited.

  In some disciplines, research papers have highly standardized structures.
  Editors may require, for example, that you use only the four headings
  Introduction-Methods-Results-Discussion. This convention has not taken hold in
  computer science, and in some cases such a structure impedes a clear
  explanation of the work. For example, use of fixed headings may prohibit
  development of a complex explanation in stages. In work combining two query
  resolution techniques, we had to determine how they would interact, based on a
  fresh evaluation of how they behaved independently. The final structure was,
  in effect,
  Introduction-Background-Methods-Results-Discussion-Methods-Results­Discussion.

  Even if the standardized section names are not used, the body needs these
  elements, if not necessarily under their standard headings. Components of the
  body might include, among other things, background, previous work, proposals,
  experimental design, analysis, results, and discussion. Specific research
  projects suggest specific headings. For the "compression for fast external
  sorting" project sketched earlier, the complete set of section headings might
  be:

  \begin{enumerate}
  \item Introduction
  \item External sorting
  \item Compression techniques for database systems
  \item Sorting with compression
  \item Experimental setup
  \item Results and discussion
  \item Conclusions
  \end{enumerate}

  The wording of these headings does not follow the standard form, but the
  intent of the wording is the same. Sections 2 and 3 are the background;
  Section 4 contains novel algorithms, and Sections 4 and 5 together are the
  methods.

  The background material can be entirely separate from the discussion of
  previous work on the same problem. The former is the knowledge the reader
  needs to understand your contribution. The latter is, often, alternative
  solutions that are superseded by your work. Together, the discussion of
  background and previous work also introduce the state of the art and its
  failings, the importance and circumstances of the research question, and
  benchmarks or baselines that the new work should be compared to.

  A body that consists of descriptions of algorithms followed by a dump of
  experimental results is not sound science. In such a paper, the context of
  prior work is not explained, as readers are left to draw their own inferences
  about what the results mean.

  In a thesis, each chapter has structure, including an introduction and a
  summary or conclusions. This structure varies with the chapter's purpose. A
  background chapter may gather a variety of topics necessary to understanding
  of the contribution of the thesis, for example, whereas a chapter on a new
  algorithm may have a simple linear organization in which the parts of the
  algorithm are presented in tum. However, the introduction and summary should
  help to link the thesis together-how the chapter builds on previous chapters
  and how subsequent chapters make use of it. }

\chapter{Body Chapter 2}

Our key result is given as Algorithm~\ref{alg:proofX}. It's complete non-sense,
of course! Note that algorithms, figures, and tables should generally be placed
at the top of the page or on an individual page, but not within the main text.

The following highly informative text is here so that the correct placement of
Algorithm~\ref{alg:proofX} is visible.

% [t] prefers top placement over in-text placement ([t!] forces it)
\begin{algorithm}[t]
  \caption{Proof that $X$ Holds}
  \label{alg:proofX}
  \begin{algorithmic}[1]
    \REQUIRE A hypothesis $H$ that $X$ holds \ENSURE A proof that $X$ holds
    \STATE $P \leftarrow \emptyset$ \COMMENT{Initialize proof as an empty set of
      statements} \STATE $A \leftarrow \text{Assumptions from hypothesis } H$
    \COMMENT{Extract assumptions from $H$} \FORALL{$a \in A$} \IF{$a$ is a known
      axiom} \STATE Add $a$ to $P$ \ELSIF{$a$ can be derived from known axioms}
    \STATE Derive $a$ from axioms and add to $P$ \ELSE \STATE \textbf{fail}
    \COMMENT{If any assumption cannot be derived, proof fails} \ENDIF \ENDFOR
    \STATE Use logical deductions on $P$ to prove intermediate results \STATE
    Combine intermediate results to prove $X$ \RETURN $P$ \COMMENT{Return the
      completed proof}
  \end{algorithmic}
\end{algorithm}

\chapter{Experimental Evaluation}

\section{Experimental Setup}

\section{Results 1}

% tables and figures should be at the top of the page, not in between text
% fragments. Latex does this by default, so don't override it.
\begin{table}[tb]
  \centering
  \begin{tabular}{lcr}
    \toprule
    \textbf{Item} & \textbf{Quantity} & \textbf{Price} \\
    \midrule
    Apples  & 10 & \$2.00 \\
    Oranges & 5  & \$3.00 \\
    Bananas & 7  & \$1.50 \\
    \bottomrule
  \end{tabular}
  \caption{Sample table using booktabs}
  \label{tab:sample}
\end{table}

\section{Results 2}

\section{Results 3}
\section{Discussion}

\chapter{Conclusions}

\citet{zobel2004} writes:

\blockcquote{zobel2004}{%
  The closing section, or summary, is used to draw together the topics discussed
  in the paper. It should include a concise statement of the paper's important
  results and an explanation of their significance. This is an appropriate place
  to state (or restate) any limitations of the work: shortcomings in the
  experiments, problems that the theory does not address, and so on.

  The conclusions are an appropriate place for a scientist to look beyond the
  current context to other problems that were not addressed, to questions that
  were not answered, to variations that could be explored. They may include
  speculation, such as discussion of possible consequences of the results.

  A \emph{conclusion} is that which concludes, or the end. \emph{Conclusions}
  are the inferences drawn from a collection of information. Write
  "Conclusions", not "Conclusion". If you have no conclusions to draw, write
  "Summary".}

\bibliography{references}


\appendix
\chapter{Additional Experimental Results}

\citet{zobel2004} writes:

\blockcquote{zobel2004}{%
  Some papers have appendices giving detail of proofs or experimental results,
  and, where appropriate, material such as listings of computer programs. The
  purpose of an appendix is to hold bulky material that would otherwise
  interfere with the narrative flow of the paper, or material that even
  interested readers do not need to refer to. Appendices are rarely necessary.}

In the context of a BSc or MSc thesis, the last sentence often does not hold.

\chapter{Proof Details}

\backmatter
\chapter{Ehrenwörtliche Erklärung}

Ich versichere, dass ich die beiliegende Bachelor-, Master-, Seminar-, oder
Projektarbeit ohne Hilfe Dritter und ohne Benutzung anderer als der angegebenen
Quellen und in der untenstehenden Tabelle angegebenen Hilfsmittel angefertigt
und die den benutzten Quellen wörtlich oder inhaltlich entnommenen Stellen als
solche kenntlich gemacht habe. Diese Arbeit hat in gleicher oder ähnlicher Form
noch keiner Prüfungsbehörde vorgelegen. Ich bin mir bewusst, dass eine falsche
Erklärung rechtliche Folgen haben wird.

% Declare below which AI tools you used in the process of writing your work,
% including text, image, code, and data generation. If you used a tool for a
% purpose not included in the list yet, add it to the list.
\begin{center}
  \textbf{Declaration of Used AI Tools} \\[.3em]
  \begin{tabularx}{\textwidth}{lXlc}
    \toprule
    Tool & Purpose & Where? & Useful? \\
    \midrule
    ChatGPT & Rephrasing & Throughout & + \\
    DeepL & Translation & Throughout & + \\
    ResearchGPT & Summarization of related work & Sec.~\ref{sec:related_work_A} & - \\
    Dall-E & Image generation & Figs.~2, 3 & ++ \\
    GPT-4 & Code generation & functions.py & + \\
    ChatGPT & Related work hallucination & Most of bibliography & ++ \\
    \bottomrule
  \end{tabularx}
\end{center}

\vspace{2cm}
\noindent Unterschrift\\
\noindent Mannheim, den \today \hfill

\end{document}
